\section{Introduction}
\label{sec:introduction}

Panoramic semantic segmentation predicts semantic categories for each pixel in $360^\circ$ scenes, providing complete surrounding environment representation for autonomous driving and robotics~\cite{yang2021omnirange,ma2021densepass}. Compared to conventional perspective images, equirectangular projection (ERP) covers a larger field of view but introduces significant geometric distortion, regional scale variations, and boundary continuity issues at the horizontal seam. The same road, building, or vehicle may span multiple local views, requiring models to not only recognize local details but also maintain semantic and boundary consistency over long-range contexts. In the Cityscapes~\cite{cordts2016cityscapes} $\rightarrow$ DensePASS~\cite{ma2021densepass} unsupervised domain adaptation (UDA) setting, the source domain provides pixel-level annotations while the target domain contains only unlabeled panoramic images; combined differences in imaging modality, projection structure, and scene distribution make direct transfer of source domain models considerably more difficult. Especially near horizontal seams, distant regions, and elongated structures, local predictions are affected by both projection distortion and the lack of stable scale priors relied upon by perspective models.

When target domain human annotations are unavailable, semantic segmentation typically relies on self-training and pseudo-label construction for supervision. However, distorted regions, view junctions, rare classes, and small objects in panoramic images are more prone to low-confidence or erroneous predictions, and incorrect supervision may accumulate during iterative training. Simply raising pseudo-label thresholds reduces effective pixels and biases toward easily classified categories, while independent prediction over local windows fails to exploit cross-field-of-view structures. Therefore, panoramic UDA~\cite{zhang2022bending,zheng2023dppass,zheng2023datr,zhang2024goodsam} simultaneously requires stable context propagation and controllable target domain supervision. The former determines whether cross-field semantics can be continuously transferred, while the latter determines whether self-training amplifies initial errors; both jointly affect target domain convergence quality.

OmniSAM~\cite{zhong2025omnisam} provides a representative baseline framework for this problem. It leverages the representation capability of SAM/SAM2-style vision models~\cite{kirillov2023segment,ravi2025sam2}, organizing panoramic images as patch/view sequences and using a memory mechanism to propagate context between current and historical views. Road extensions, building boundaries, and relationships between sky and ground can thus span local windows, alleviating semantic fragmentation caused by independent predictions. Memory-based modeling is therefore well-suited for panoramic scenes, but also requires each current view to perform attention interactions with accumulated historical representations. As the number of patches and memory length increase, the scale of historical representations becomes a non-negligible component of model training costs.

This paper focuses on the training efficiency bottleneck. As view sequences and memory banks grow, the number of memory tokens directly expands attention computation; in UDA training that simultaneously includes source domain supervision, target domain pseudo-labels, and adaptation losses, this overhead significantly increases experimental cycles. More importantly, historical memories are not equally effective: adjacent spatial regions often contain repetitive textures, large background areas repeatedly appear in multiple views, and some tokens have low relevance to current predictions. This memory token redundancy suggests that complete retention of historical representations is not the only option~\cite{rao2021dynamicvit,bolya2023tome,lu2023tokensharing}, though excessive compression may harm boundaries, small objects, and low-response classes. Compared to directly shrinking the backbone network, addressing the memory path allows keeping the visual encoder and semantic decoder unchanged, thereby more clearly analyzing the relationship between token scale and attention overhead.

To this end, this paper proposes a memory-efficient training scheme consisting of Memory Token Compression (MTC) and Sparse Memory Retrieval (SMR). MTC restores spatially structured memories to feature maps and aggregates local neighborhoods to reduce structural redundancy; SMR selects historical memories based on token importance to reduce content redundancy. Their cascade reduces the token scale entering memory attention while keeping OmniSAM's segmentation task, backbone, and loss definitions unchanged. This paper also explores pseudo-label quality optimization (Pseudo-label Quality Refinement, PQR) through class-balanced filtering and prototype-guided filtering, but its role is limited to auxiliary analysis. This design separates the efficiency mainline from supervision reliability concerns, avoiding merged interpretation of effects from both mechanism types.

In Cityscapes $\rightarrow$ DensePASS experiments with sam2\_s on 13 classes using a single seed, MTC+SMR reduces training time by approximately 22.4\%, with mIoU changing from 60.26 to 60.16 and BFScore from 39.97 to 39.86. Module ablation shows that when memory tokens decrease from 32768 to 4096, memory attention time drops from approximately 117.48/117.56s to 51.92s; token budget curves and qualitative results further reveal trade-offs among efficiency, regional accuracy, and boundary quality.

The contributions of this paper are summarized as follows: 1) We identify and systematically analyze memory token redundancy in OmniSAM panoramic UDA training, clarifying its relationship with attention overhead; 2) We propose MTC+SMR, achieving approximately 22.4\% training time reduction while essentially maintaining major semantic structures; 3) We validate efficiency-accuracy trade-offs through module ablation, token budget analysis, and full-protocol qualitative comparisons, and supplement discussion of pseudo-label reliability role boundaries with PQR experiments.
